\documentclass[12pt,a4paper]{article}

\usepackage[utf8]{inputenc}
\usepackage[T1]{fontenc}
\usepackage[brazil]{babel}
\usepackage{geometry}
\usepackage{graphicx}
\usepackage{xcolor}
\usepackage{hyperref}
\usepackage{longtable}
\usepackage{array}
\usepackage{enumitem}
\usepackage{titlesec}
\usepackage{fancyhdr}

\geometry{a4paper, margin=2.5cm}
\setlength{\parindent}{0pt}
\setlength{\parskip}{0.8em}
\setlist[itemize]{leftmargin=1.2cm}
\hypersetup{
    colorlinks=true,
    linkcolor=blue,
    urlcolor=blue
}

\definecolor{prototipeAzul}{HTML}{0397DD}
\definecolor{prototipeRoyal}{HTML}{1D58AA}
\definecolor{prototipeMarinho}{HTML}{060D5B}
\definecolor{prototipeCreme}{HTML}{FDFBDB}

\titleformat{\section}
  {\normalfont\Large\bfseries\color{prototipeMarinho}}
  {\thesection}{1em}{}

\titleformat{\subsection}
  {\normalfont\large\bfseries\color{prototipeRoyal}}
  {\thesubsection}{1em}{}

\pagestyle{fancy}
\fancyhf{}
\lhead{Chave 26}
\rhead{Relatório do Projeto}
\cfoot{\thepage}

\begin{document}

\begin{titlepage}
    \centering
    \vspace*{2cm}

    {\Huge\bfseries\color{prototipeMarinho} Relatório do Projeto Chave 26\par}
    \vspace{0.7cm}
    {\Large Aplicativo para controle simples da Sala 26 da Prototipe\par}
    \vspace{1.5cm}

    {\large Relatório em linguagem simples, com resumo do funcionamento, tecnologias, decisões, limitações e estado atual do projeto.\par}

    \vfill

    {\large Projeto: \textbf{Chave 26}\par}
    {\large Plataforma principal: \textbf{Android / celular}\par}
    {\large Data do relatório: \textbf{29 de maio de 2026}\par}

    \vspace*{1.5cm}
\end{titlepage}

\tableofcontents
\newpage

\section{Resumo executivo}

O \textbf{Chave 26} é um aplicativo criado para resolver um problema simples do dia a dia da Prototipe: saber rapidamente se a sala 26 está aberta ou fechada, onde está a chave e quem foi a última pessoa responsável pela movimentação.

Hoje, esse tipo de informação costuma depender de mensagens em grupo, como ``sala aberta'' ou ``chave comigo''. O problema é que, com o tempo, essas mensagens ficam espalhadas e pode ser difícil saber qual é a situação mais recente. O aplicativo centraliza essa informação em uma tela simples.

O projeto foi construído como um \textbf{MVP}, ou seja, uma primeira versão funcional para validar a ideia. Ele não tenta ser um sistema completo de produção. A prioridade foi entregar algo funcional, claro, bonito, fácil de apresentar e com regras bem definidas.

\section{Objetivo do projeto}

O objetivo principal do Chave 26 é permitir que uma pessoa abra o aplicativo e entenda, em poucos segundos:

\begin{itemize}
    \item se a sala 26 está aberta ou fechada;
    \item onde está a chave;
    \item se a chave está com alguém, na sala ou guardada em algum local;
    \item quem fez a última alteração;
    \item quando essa alteração aconteceu;
    \item quais foram as movimentações recentes.
\end{itemize}

Em linguagem simples, o aplicativo funciona como um painel de situação da sala. Ele mostra o estado atual e guarda um histórico das ações importantes.

\section{Público-alvo}

O aplicativo foi pensado para membros e trainees da Prototipe que usam a sala 26 ou precisam saber se ela está disponível.

Como esta é uma versão de demonstração, o app usa perfis locais pré-definidos, em vez de login real. Isso permite simular pessoas diferentes usando o mesmo celular durante uma apresentação.

\section{O que o aplicativo faz}

\subsection{Tela principal}

A tela principal mostra o nome do app, o perfil selecionado, a situação da sala e as ações que podem ser feitas naquele momento. A interface foi organizada para dar prioridade ao que importa primeiro: saber se a sala está aberta ou fechada e onde está a chave.

As principais informações exibidas são:

\begin{itemize}
    \item status da sala: aberta ou fechada;
    \item pessoa responsável, quando houver;
    \item localização da chave;
    \item horário da última atualização;
    \item histórico recente.
\end{itemize}

\subsection{Perfis locais}

O aplicativo possui quatro perfis locais para demonstração:

\begin{itemize}
    \item Lucas;
    \item Clara;
    \item Amanda;
    \item Vitor.
\end{itemize}

Esses perfis não são contas reais. Eles servem para identificar quem está realizando cada ação durante o protótipo. A troca de perfil acontece pela própria tela do app e a escolha fica salva no aparelho.

\subsection{Ações da chave e da sala}

O aplicativo permite registrar ações ligadas à sala e à chave, como:

\begin{itemize}
    \item pegar a chave na portaria ou em outro local onde ela esteja guardada;
    \item abrir a sala;
    \item fechar a sala;
    \item guardar a chave em um destino escolhido;
    \item passar a chave para outra pessoa;
    \item registrar automaticamente a movimentação no histórico.
\end{itemize}

As ações aparecem conforme a situação atual. Por exemplo, se a chave está com outra pessoa, o app não oferece uma ação que o usuário atual não pode fazer. Isso reduz erro e deixa a experiência mais simples.

\subsection{Destinos da chave}

Além da portaria, o aplicativo permite que a chave seja guardada em outros locais. Existem destinos fixos e destinos personalizados.

Destinos fixos atuais:

\begin{itemize}
    \item Portaria;
    \item Maker Space.
\end{itemize}

Também é possível cadastrar, renomear e remover destinos personalizados no próprio app. Esses destinos ficam salvos localmente no aparelho.

\subsection{Histórico}

Toda ação importante gera um item no histórico. O histórico mostra:

\begin{itemize}
    \item quem fez a ação;
    \item qual foi a ação;
    \item data e horário da movimentação.
\end{itemize}

O histórico aparece em ordem do mais recente para o mais antigo. Também existe uma aba própria para consultar o histórico completo, incluindo filtro por data.

\section{Tecnologias utilizadas}

\begin{longtable}{|p{4cm}|p{9cm}|}
\hline
\textbf{Tecnologia} & \textbf{Uso no projeto} \\
\hline
Flutter & Framework usado para construir o aplicativo mobile. Ele permite criar interfaces visuais para celular usando uma única base de código. \\
\hline
Dart & Linguagem de programação usada pelo Flutter. Todo o app foi escrito em Dart. \\
\hline
Material 3 & Base visual usada para botões, menus, diálogos, navegação e componentes de interface. \\
\hline
Riverpod & Biblioteca usada para organizar o estado do app. No projeto, ela é usada de forma simples, sem arquitetura excessiva. \\
\hline
SharedPreferences & Tecnologia usada para salvar dados pequenos diretamente no aparelho, como perfil selecionado, situação da sala, histórico e destinos customizados. \\
\hline
JSON & Formato usado para transformar os dados da sala em texto e salvar no armazenamento local. \\
\hline
Flutter Test & Ferramenta usada para criar testes automatizados de regras e de interface. \\
\hline
Android SDK & Conjunto de ferramentas usado para gerar e testar o aplicativo Android. \\
\hline
OpenJDK 17 & Ferramenta necessária para o processo de build Android. \\
\hline
\end{longtable}

\section{Organização técnica do projeto}

O projeto foi organizado de forma simples, separando as partes principais:

\begin{itemize}
    \item \textbf{Domínio}: contém as regras da sala, os estados possíveis e as ações permitidas.
    \item \textbf{Dados}: contém a parte de salvar e carregar informações no aparelho.
    \item \textbf{Interface}: fica concentrada na tela principal do app.
    \item \textbf{Testes}: validam regras importantes e comportamentos da tela.
    \item \textbf{Documentação}: registra decisões, progresso, funcionalidades e limitações.
\end{itemize}

Os nomes ligados ao assunto do app foram escritos em português, como \texttt{EstadoDaSala}, \texttt{LocalizacaoDaChave}, \texttt{SituacaoDaSala}, \texttt{EventoHistorico} e \texttt{ControladorDaSala}. Isso facilita a explicação do projeto para pessoas que estão aprendendo ou avaliando o código.

\section{Principais regras de funcionamento}

O aplicativo possui regras para evitar registros incoerentes. Algumas delas são:

\begin{itemize}
    \item a sala só pode ser aberta por quem está com a chave;
    \item não é possível abrir a sala se ela já está aberta;
    \item não é possível fechar a sala se ela já está fechada;
    \item não é possível transferir a chave para a própria pessoa;
    \item não é possível registrar movimentação sem perfil selecionado;
    \item nomes vazios são bloqueados nas ações importantes;
    \item guardar a chave exige informar um destino;
    \item toda ação válida registra um evento no histórico;
    \item ações inválidas não alteram a situação salva.
\end{itemize}

Essas regras ficam protegidas no código principal de domínio, não apenas na tela. Isso é importante porque a interface ajuda o usuário, mas a regra real precisa estar protegida na lógica do sistema.

\section{Persistência dos dados}

O app funciona de forma local e offline. Isso significa que os dados ficam salvos no próprio aparelho.

São salvos localmente:

\begin{itemize}
    \item situação atual da sala;
    \item localização atual da chave;
    \item histórico de movimentações;
    \item perfil local selecionado;
    \item destinos personalizados da chave.
\end{itemize}

A situação da sala e o histórico são salvos juntos em um único JSON. Para este protótipo, isso torna a implementação mais simples e reduz o risco de o estado atual ficar separado do histórico.

\section{Identidade visual}

O projeto usa a identidade visual da Prototipe como referência. A documentação registra uma paleta com tons de azul, ciano, índigo, marinho, branco e creme.

Também foram usados assets locais, incluindo:

\begin{itemize}
    \item logo;
    \item mascote da Prototipe;
    \item ícone do app Chave 26.
\end{itemize}

O aplicativo Android também recebeu personalização nativa:

\begin{itemize}
    \item nome exibido como ``Chave 26'';
    \item ícone personalizado no launcher;
    \item splash screen nativa com logo;
    \item configuração específica para Android 12 ou superior.
\end{itemize}

\section{Testes e qualidade}

O projeto possui testes automatizados cobrindo regras de negócio, persistência e interface.

Na verificação realizada para este relatório:

\begin{itemize}
    \item \texttt{flutter test}: todos os testes passaram;
    \item \texttt{flutter analyze}: nenhum problema encontrado.
\end{itemize}

Os testes cobrem pontos como:

\begin{itemize}
    \item criação e comparação dos modelos da sala;
    \item abertura e fechamento da sala;
    \item transferência da chave;
    \item bloqueio de ações inválidas;
    \item salvamento e carregamento local;
    \item restauração dos dados ao abrir o app;
    \item histórico vazio e histórico com itens;
    \item filtro de histórico por data;
    \item destinos personalizados;
    \item troca de perfil;
    \item visibilidade das ações conforme a situação atual.
\end{itemize}

\section{Documentação existente}

O projeto possui documentação contínua, distribuída em arquivos específicos:

\begin{itemize}
    \item \texttt{README.md}: visão geral do projeto e comandos principais.
    \item \texttt{CONTEXTO.md}: problema, objetivo, público e decisões gerais.
    \item \texttt{BACKLOG.md}: planejamento, marcos, progresso e pendências.
    \item \texttt{TRADEOFFS.md}: escolhas conscientes e limitações aceitas.
    \item \texttt{documentos/decisoes-tecnicas.md}: stack, arquitetura e decisões técnicas.
    \item \texttt{documentos/plano-inicial.md}: plano original do MVP.
    \item \texttt{documentos/identidade-visual-prototipe.md}: cores e referências visuais.
    \item \texttt{documentos/funcionalidades/}: registro das funcionalidades implementadas por etapa.
\end{itemize}

\section{Fluxogramas do projeto}

Os fluxogramas abaixo registram visualmente a evolução dos marcos do projeto e ajudam a entender, de forma mais direta, como as etapas foram organizadas.

\begin{figure}[htbp]
    \centering
    \includegraphics[width=\textwidth]{marco1.png}
    \caption{Fluxograma do Marco 1.}
\end{figure}

\begin{figure}[htbp]
    \centering
    \includegraphics[width=\textwidth]{marco2.png}
    \caption{Fluxograma do Marco 2.}
\end{figure}

\begin{figure}[htbp]
    \centering
    \includegraphics[width=\textwidth]{marco3.png}
    \caption{Fluxograma do Marco 3.}
\end{figure}

\begin{figure}[htbp]
    \centering
    \includegraphics[width=\textwidth]{marco4.png}
    \caption{Fluxograma do Marco 4.}
\end{figure}

\begin{figure}[htbp]
    \centering
    \includegraphics[width=0.62\textwidth]{marco5.png}
    \caption{Fluxograma do Marco 5.}
\end{figure}

\clearpage

\section{O que ficou dentro do MVP}

Entraram no MVP:

\begin{itemize}
    \item app mobile com foco em Android;
    \item funcionamento local e offline;
    \item perfis locais pré-definidos;
    \item estado da sala;
    \item localização da chave;
    \item ações principais da chave e da sala;
    \item transferência da chave;
    \item histórico simples;
    \item destinos fixos e customizados;
    \item persistência local;
    \item confirmações para ações importantes;
    \item interface principal organizada;
    \item ícone e abertura nativa Android personalizados;
    \item testes automatizados;
    \item documentação do projeto.
\end{itemize}

\section{O que ficou fora do MVP}

Ficaram fora desta primeira versão:

\begin{itemize}
    \item backend;
    \item login real;
    \item autenticação por senha;
    \item sincronização entre vários celulares;
    \item notificações;
    \item integração com grupo de mensagens;
    \item painel web;
    \item modo escuro;
    \item auditoria oficial contra alterações locais.
\end{itemize}

Esses itens não foram esquecidos. Eles ficaram fora porque aumentariam o tamanho e o risco do projeto. Para um protótipo de trainee, a decisão foi priorizar um app simples, funcional e demonstrável.

\section{Limitações atuais}

As principais limitações do projeto são:

\begin{itemize}
    \item os dados ficam apenas no aparelho onde o app está instalado;
    \item outro celular não vê automaticamente a mesma situação da sala;
    \item não há login real, então o perfil selecionado não prova identidade;
    \item qualquer pessoa com acesso ao aparelho pode trocar o perfil;
    \item o histórico local não deve ser tratado como auditoria oficial;
    \item ao desinstalar o app, os dados podem ser perdidos;
    \item \texttt{SharedPreferences} é suficiente para poucos dados, mas não é a melhor escolha para histórico grande;
    \item não há busca avançada, relatórios ou sincronização;
    \item destinos personalizados são locais ao aparelho;
    \item a confirmação visual para apagar destino customizado ainda pode ser melhorada;
    \item o tema visual ainda está concentrado em \texttt{lib/main.dart} e pode ser separado em uma pasta própria futuramente.
\end{itemize}

\section{Possíveis evoluções futuras}

Se o projeto virar um sistema usado oficialmente, os próximos passos recomendados seriam:

\begin{itemize}
    \item criar login real para membros;
    \item sincronizar a situação da sala entre vários celulares;
    \item usar backend com banco de dados central;
    \item adicionar notificações;
    \item criar auditoria confiável das movimentações;
    \item permitir permissões por tipo de usuário;
    \item criar busca e filtros mais completos no histórico;
    \item gerar relatórios de uso;
    \item transformar destinos em configuração administrável;
    \item revisar a persistência local e migrar para banco mais robusto, se necessário.
\end{itemize}

Tecnologias possíveis para uma versão mais completa incluem Firebase, Supabase, API própria, SQLite, Drift, Hive ou Isar. A escolha dependeria do tamanho real do uso e da necessidade de sincronização.

\section{Conclusão}

O Chave 26 entrega uma solução simples para um problema real: saber a situação da sala 26 e da chave sem depender de mensagens espalhadas em grupo.

O projeto foi construído como um protótipo funcional, com foco em Android, funcionamento offline, interface clara, regras de negócio testadas, persistência local e documentação consistente.

Ele ainda não é uma solução oficial multiusuário, porque não possui login real, backend ou sincronização. Mesmo assim, como MVP, cumpre bem o papel de validar a ideia, demonstrar o fluxo principal e mostrar uma base técnica organizada para futuras evoluções.

\end{document}
